\section{Related Work}
\label{sec:related}

The tools closest to \sys fall short along three primary dimensions: operator-level timing
without counters, kernel-level counter collection without operator identity, and
compilation-level optimization without counter-verified feedback. Adjacent areas---fused kernel
libraries and profile-guided optimization---are surveyed below for completeness.

\paragraph{Operator-level profilers.}
\texttt{torch.profiler} and its Kineto backend \cite{kineto2021} record wall-clock durations and
Chrome traces at the \texttt{aten::} operator granularity~\cite{pytorch2019} and expose CUDA activity correlation
via CUPTI External IDs---giving kernel-level \emph{timing} attribution but no access to hardware
performance counters. Practitioners can identify \emph{which} operator is slow, but cannot
determine \emph{which hardware resource} it is contending for---the distinction that determines
what optimization is warranted. Without hardware
counter data at the operator level, optimization choices remain heuristic.
NVIDIA DLProf~\cite{dlprof} is the closest NVIDIA-specific predecessor: it uses NVTX markers
to correlate CPU and GPU time with model operations, reports at operation, layer, and kernel
granularity, and includes Tensor Core usage diagnostics. DLProf remains timing-centric:
its GPU-side data comes from \nsys, and Tensor Core utilization is inferred from kernel names.
\sys adds the missing counter-attribution layer. It joins application-mode \ncu hardware
counter rows to \nsys kernel records through invocation-order matching, then assigns those
counter-enriched kernel records to PyTorch operators, including Inductor-era fused kernels.
This produces per-operator hardware metrics such as measured Tensor Core activity,
occupancy, memory throughput, register pressure, and cache behavior rather than only
operation timing or kernel-name-derived Tensor Core eligibility.

\paragraph{Kernel-level profilers.}
NVIDIA Nsight Systems (\nsys) \cite{nsightsystems} captures the full CUDA kernel timeline with
NVTX annotations at high throughput, but produces no operator-level view: every kernel appears by
its internal CUDA or Triton name (\eg \texttt{cutlass\_80\_simt\_sgemm\_128x32\_tn}), with no
indication of which PyTorch operator dispatched it. NVIDIA Nsight Compute (\ncu)
\cite{nsightcompute} collects per-kernel hardware counters at full precision---occupancy, DRAM
bytes, Tensor Core utilization, cache hit rates---but its output is keyed on kernel names and
invocation indices in the replay timeline, again with no awareness of the operator graph above.
NVIDIA DCGM and Dynolog \cite{dynolog2023} operate at cluster monitoring granularity and do not
perform per-kernel or per-operator attribution. Taken together, these tools collectively address
hardware counter collection and kernel-timeline capture but leave kernel-to-operator attribution entirely unresolved.

\paragraph{Deep learning compilers and optimizers.}
TorchDynamo and TorchInductor \cite{pytorch2024} compile PyTorch programs via FX graph capture~\cite{torchfx2022} to Triton kernels
using heuristic fusion rules (pointwise chaining, matmul epilogue fusion), but fusion decisions
are made without hardware counter feedback: whether a given fusion improved Tensor Core
utilization is not a question the standard Inductor toolchain can answer programmatically. TVM \cite{tvm2018} and Halide \cite{halide2013} optimize tensor programs via
measurement-based auto-tuning and schedule/algorithm separation respectively; both optimize
\emph{latency} directly rather than using hardware counters to identify the bottleneck class
driving that latency. ONNX Runtime\footnote{\url{https://onnxruntime.ai}} and NVIDIA
TensorRT\footnote{\url{https://developer.nvidia.com/tensorrt}} perform layer fusion and
quantization behind black-box interfaces that cannot be extended for custom architectures or
augmented with attribution data.

\paragraph{Fused kernel libraries.}
FlashAttention \cite{flashattention2022,flashattention2024} and xFormers \cite{xformers2022}
provide hand-optimized fused attention kernels that eliminate DRAM round-trips by keeping
intermediate tensors in on-chip SRAM. These represent the class of optimization that
attributed hardware counter evidence should motivate rather than apply unconditionally; what
is absent is a profiling layer that identifies, from measured per-operator counters, which
variant is warranted for a given workload.

\paragraph{Profile-guided optimization.}
Classic PGO in C/C++ compilers (GCC, LLVM) collects branch frequency profiles to guide inlining
and layout decisions \cite{autofdo2016}. Applying the same principle at the GPU operator level requires hardware counter data attributed
to specific operators. To our knowledge, \sys is the first published system to combine
\nsys kernel traces, application-mode \ncu hardware counters, and PyTorch operator
attribution through invocation-order matching, and to use the resulting attributed profile
to drive hardware-informed FX graph optimization.
